\documentclass[11pt,letterpaper]{article}
\usepackage[T1]{fontenc}
\usepackage[utf8]{inputenc}
\usepackage[spanish,es-noshorthands]{babel}
\usepackage{lmodern}
\usepackage{microtype}
\usepackage[margin=2.5cm]{geometry}
\usepackage{amsmath}
\usepackage{booktabs}
\usepackage{enumitem}
\usepackage{array}
\pagestyle{empty}
\setlength{\parindent}{0pt}
\setlength{\parskip}{6pt}

\begin{document}

{\large\textbf{Econometría II · Preparación PEP 1}}\hfill 1 de noviembre de 2026

\vspace{2pt}\hrule\vspace{10pt}

{\Large\textbf{Ejercicio 3. DiD $2\times 2$}}\hfill Tiempo: 20 min

\textit{Ejercicio de práctica con datos ficticios; no proviene de una prueba del profesor.}

\vspace{4pt}
Se quiere evaluar el traspaso de escuelas municipales a los Servicios Locales de Educación Pública (SLEP). SIMCE promedio por grupo:

\begin{center}
\begin{tabular}{l>{\raggedleft\arraybackslash}p{2cm}>{\raggedleft\arraybackslash}p{2cm}}
\toprule
 & \textbf{Antes} & \textbf{Después} \\
\midrule
SLEP (traspasadas)       & 240 & 252 \\
Municipal (siguen igual) & 255 & 259 \\
\bottomrule
\end{tabular}
\end{center}

\[
\text{SIMCE}=\beta_0+\beta_1\,\text{SLEP}+\beta_2\,\text{Post}+\beta_3\,(\text{SLEP}\times\text{Post})+u
\]

El error estándar de $\hat{\beta}_3$ es 3,2. Supón una muestra grande.

\begin{enumerate}[label=\textbf{\alph*)},leftmargin=*,itemsep=8pt]
  \item Calcula el DiD desde las medias e interprétalo.
  \item Calcula las dos comparaciones ingenuas (antes y después solo en SLEP; SLEP contra municipal solo después). ¿Qué mezcla cada una?
  \item Da los cuatro $\hat{\beta}$ y di qué mide cada uno.
  \item Calcula el estadístico $t$ de $\hat{\beta}_3$ y concluye al 5\,\%.
  \item ¿Qué exige aquí el supuesto de tendencias paralelas? ¿Cuál sería el SIMCE de las escuelas SLEP en el periodo «después» sin el traspaso?
\end{enumerate}

\vspace{6pt}
\textit{Cada resultado se entrega con su frase de interpretación.}

\end{document}
